% =============================================================================
% SECTION 3 — METHOD
% (Updated for the new minimal three-panel Figure~\ref{fig:pdg_score} produced
%  by output_new/plot_pdg_score.py.  All numerical detail formerly carried by
%  the figure caption now lives in the body paragraphs and the worked-example
%  appendix.)
% =============================================================================
% Forward references used:
%   Section~\ref{sec:exp:ablation}
%   Section~\ref{sec:method:cot}
%   Appendix~\ref{app:algo}
%   Appendix~\ref{app:algo:worked_example}
% =============================================================================

\section{Method}
\label{sec:method}

\subsection{Motivation: Function Calls as Agent-Like Structures}
\label{sec:method:motivation}

Formally, at step $t$ a coding agent maintains a history $h_t$ and samples
an action $a_t \sim \pi(a_t \mid h_t)$; the environment yields an
observation $o_{t+1} \sim p(o_{t+1} \mid h_t, a_t)$; and the agent
continues conditioned on the full trace,
$\pi(\,\cdot \mid h_t, a_t, o_{t+1})$. Under the function-call/agent-step
correspondence introduced in Section~\ref{sec:intro}
(Figure~\ref{fig:teaser}, left), $h_t$ aligns with the pre-call context,
$a_t$ with the call, $o_{t+1}$ with the return, and the continuation with
downstream usage of the returned value.

This correspondence motivates our central design: extract an agentic
conditioning signal from the structure already present in source code via
the fill-in-the-middle (FIM) objective. The correspondence is
\emph{structural} rather than literal---FIM conditions on a given suffix
while agent inference generates one---but training a model to reason
bidirectionally about function-level dependencies should induce a
representation useful when the ``suffix'' is generated rather than
observed. Standard random-span
FIM~\citep{qwencoder, deepseekcoder, starcoder} captures this
structure only incidentally; we instead construct a \emph{function-aware}
variant that selects masking targets based on program structure and
contextual predictability.

\subsection{Data Collection and Decontamination}
\label{sec:method:data}

We curate a corpus of $990$ Python repositories from GitHub. Starting from
$\sim\!2{,}000$ candidates retrieved by combining a star-count threshold with
ten topic queries (Appendix~\ref{app:data}), we apply manual quality filtering
and remove every repository whose origin overlaps with the source repositories
of SWE-Bench~\citep{swebench} (verified by repository name and any known
fork). To eliminate the possibility of test-time leakage from training-data
commits, we restrict each repository to commits whose timestamp precedes the
earliest base-commit used in SWE-Bench-Verified and SWE-Bench-Lite. The full
repository list, license inventory, and commit hashes will be released to
support reproduction.

% -----------------------------------------------------------------------------
% Figure 3 (categories) and Table 1 (corpus stats) placed side-by-side
% -----------------------------------------------------------------------------
\begin{figure}[t]
\centering
\begin{minipage}[t]{0.50\linewidth}
  \centering
  \includegraphics[width=\linewidth]{figs/dataset_categories.pdf}
  \caption{Distribution of the 990 source repositories across
  ten topic categories. The corpus is dominated by reference
  implementations, scientific computing, and small frameworks; compilers
  and networking/security tails remain by design to maintain coverage
  diversity.}
  \label{fig:dataset_categories}
\end{minipage}\hfill
\begin{minipage}[t]{0.46\linewidth}
  \vspace{0pt}
  \centering
  \caption{Mid-training corpus statistics after decontamination
  and quality filtering. Token counts are reported with the Qwen2.5-Coder
  tokenizer.}
  \label{tab:dataset}
  \scriptsize
  \setlength{\tabcolsep}{4pt}
  \begin{tabular}{lr}
  \toprule
  Property & Value \\
  \midrule
  Source GitHub repositories                       & $990$ \\
  Topic categories                                 & $10$ \\
  Filtered self-contained Python files             & $\approx 78\mathrm{K}$ \\
  Total lines of Python code                       & $\approx 33\mathrm{M}$ \\
  Mean LoC per file                                & \textcolor{red}{$\approx 425$} \\
  Mid-training token budget                        & $\approx 3\mathrm{B}$ \\
  \midrule
  Single-function FIM targets                      & \textcolor{red}{$\approx 412\mathrm{K}$} \\
  Multi-function targets ($k{=}2$)                 & \textcolor{red}{$\approx 96\mathrm{K}$} \\
  Multi-function targets ($k{=}3$)                 & \textcolor{red}{$\approx 31\mathrm{K}$} \\
  Mean target LoC                                  & \textcolor{red}{$\approx 38$} \\
  Targets with Gemini-3 CoT                        & \textcolor{red}{$100\%$} \\
  \midrule
  SWE-Bench source-repo overlap                    & $0$ \\
  Commit cutoff                                    & before earliest SWE-Bench base \\
  License inventory                                & permissive only \\
  \bottomrule
  \end{tabular}
\end{minipage}
\end{figure}

\subsection{Function-Aware FIM Target Selection}
\label{sec:method:selection}

Given the file source of each repository, our pipeline produces a set of
mask targets together with the corresponding masked file. Each target is
a single function or a connected group of $2$--$3$ functions; the latter
variant is studied in Section~\ref{sec:exp:ablation}. The pipeline has
four stages: dependency-graph construction, complexity scoring,
inferability scoring, and threshold-based selection. We illustrate all
four stages in Figure~\ref{fig:pdg_score} on a tiny running example---a
\texttt{Calculator} class with two top-level helpers
\texttt{add}/\texttt{is\_int}. The full single-function selection
algorithm is in Appendix~\ref{app:algo:single}; a line-by-line numerical
walkthrough of the running example---reproducing every quantity behind
Figure~\ref{fig:pdg_score}---is in
Appendix~\ref{app:algo:worked_example}.

% -----------------------------------------------------------------------------
% Figure: PDG construction and H/I score schematic
% -----------------------------------------------------------------------------
\begin{figure}[t]
\centering
\includegraphics[width=\linewidth]{figs/pdg_score.pdf}
\caption{Function-aware FIM target selection on a tiny calculator
example. \textbf{(a)} The program dependency graph parsed from the
file. \textbf{(b)} Each function receives a complexity score $\hat{H}$
(three structural components) and an inferability score $\hat{I}$
(five contextual components), combined into a harmonic-mean-like
score. \textbf{(c)} Only functions in the rose-shaded region
$\mathrm{FIM}\!\geq\!\tau$ become mask targets.}
\label{fig:pdg_score}
\end{figure}

\subsubsection{Program Dependency Graph}
\label{sec:method:pdg}

For each source file we parse its AST and extract the set $\mathcal{V}$
of function nodes (top-level functions and class methods, identified by
qualified names). We construct two edge sets: \emph{call edges}
$\mathcal{E}_{\mathrm{call}}$ between caller and callee, and
\emph{sibling edges} $\mathcal{E}_{\mathrm{sib}}$ between methods of the
same class (capturing intra-class coupling that does not manifest as
direct calls but carries information through shared instance state).
Call resolution handles common Python idioms (direct invocations, class
instantiation, \texttt{self}/\texttt{cls} method calls) with a
short-name fallback against the qualified-name index; details are in
Appendix~\ref{app:algo:pdg}. In
Figure~\ref{fig:pdg_score}(a), \texttt{Calculator.total} reaches the
top-level helpers via $\mathcal{E}_{\mathrm{call}}$, while
$\mathcal{E}_{\mathrm{sib}}$ connects the four \texttt{Calculator}
methods that share \texttt{self.history}.

\subsubsection{Complexity Score $\hat{H}$}
\label{sec:method:complexity}

For each $v \in \mathcal{V}$ we define
%
\begin{equation}
  \hat{H}(v) \;=\; w_{\ell}\,\phi\!\left(\mathrm{LoC}(v),\,c_{\ell}\right)
            + w_{c}\,\phi\!\left(\mathrm{CC}(v),\,c_{c}\right)
            + w_{d}\,\phi\!\left(\mathrm{D}(v),\,c_{d}\right),
  \label{eq:complexity}
\end{equation}
%
where $\mathrm{LoC}(v)$ is lines of code, $\mathrm{CC}(v)$ is McCabe
cyclomatic complexity, $\mathrm{D}(v)$ is the maximum nesting depth of
control-flow constructs, and $\phi(x, c) = \min(x / c,\, 2)$ normalizes
each quantity by a soft cap of $2$ (allowing genuinely complex functions
to score above the median without letting outliers dominate). The
default weights and caps are
$(w_{\ell},w_{c},w_{d})\!=\!(0.4,0.4,0.2)$ and
$(c_{\ell},c_{c},c_{d})\!=\!(50,10,5)$
(Appendix~\ref{app:algo:hyperparams}). On the running example,
\texttt{Calculator.total} has $\mathrm{LoC}\!=\!10$,
$\mathrm{CC}\!=\!5$, and $\mathrm{D}\!=\!3$, giving
$\hat{H}\!=\!0.08\!+\!0.20\!+\!0.12\!=\!0.40$;
the upper bar of Figure~\ref{fig:pdg_score}(b) visualizes this as a
three-segment stack. The line-by-line derivation of these counts is in
Appendix~\ref{app:algo:worked_example}.

\subsubsection{Inferability Score $\hat{I}$}
\label{sec:method:inferability}

A target should be \emph{recoverable} from the surrounding context.
$\hat{I}(v)$ aggregates five context-derived signals that approximate the
mutual information between $v$'s body and the rest of the file:
%
\begin{equation}
  \hat{I}(v) \;=\; \alpha\,C_{\mathrm{caller}}(v)
              + \beta\,C_{\mathrm{callee}}(v)
              + \gamma\,C_{\mathrm{sig}}(v)
              + \delta\,C_{\mathrm{doc}}(v)
              + \varepsilon\,C_{\mathrm{class}}(v).
  \label{eq:inferability}
\end{equation}
%
$C_{\mathrm{caller}}$ scores call-site argument specificity,
$C_{\mathrm{callee}}$ counts intra-file functions called by $v$,
$C_{\mathrm{sig}}$ aggregates type annotations and name descriptiveness,
$C_{\mathrm{doc}}$ indicates docstring presence, and
$C_{\mathrm{class}}$ counts in-class siblings. Default mixing weights
are $(\alpha,\beta,\gamma,\delta,\varepsilon)\!=\!(0.30,0.25,0.20,
0.10,0.15)$ and full component formulas are in
Appendix~\ref{app:algo:inferability}. Each component is a hand-designed
proxy rather than a learned predictability score: a learned proxy would
couple selection to a particular reference model and complicate the
cross-base-model generalization analysis (Section~\ref{sec:exp:main}).
For \texttt{Calculator.total} the five terms evaluate to $0.06$ (one
caller, \texttt{mean}), $0.13$ (two intra-file callees,
\texttt{is\_int}/\texttt{add}), $0.10$ (return-type annotation plus a
single-token name), $0.05$ (docstring present), and $0.14$ (three
siblings together with the existing \texttt{\_\_init\_\_} bonus),
summing to $\hat{I}\!=\!0.48$. The lower bar of
Figure~\ref{fig:pdg_score}(b) visualises these proportions; the
class-context term dominates because \texttt{total} is embedded in a
four-method class. A component-by-component derivation is in
Appendix~\ref{app:algo:worked_example}.

\subsubsection{Single-Function Score}
\label{sec:method:single}

The final score combines $\hat{H}$ and $\hat{I}$ in a harmonic-mean-like
form that requires both to be simultaneously high, multiplied by a
one-sided difficulty penalty $\rho(\Delta(v))$:
%
\begin{equation}
  \mathrm{FIM}(v) \;=\;
  \frac{\hat{H}(v)\,\hat{I}(v)}{\hat{H}(v) + \hat{I}(v) + \epsilon}
  \cdot \rho\!\left(\Delta(v)\right),
  \label{eq:fim_single}
\end{equation}
%
where $\epsilon$ is a small constant. The harmonic-mean form forces both
$\hat{H}$ and $\hat{I}$ to be large simultaneously, penalizing imbalance;
the penalty $\rho$ down-weights targets that remain hard even given full
context (which would otherwise be unlearnable noise). Full forms of
$\Delta$ and $\rho$, the hard filters on length, dunder methods, and
minimum scores, and the threshold $\tau\!=\!0.20$ used throughout this
work are documented in Appendix~\ref{app:algo:penalty}.
Figure~\ref{fig:pdg_score}(c) plots all candidate functions in the
$(\hat{H},\hat{I})$ plane: the four short helpers fall below the
LoC hard filter and are removed before scoring, while
\texttt{Calculator.total} attains
$\mathrm{FIM}\!=\!0.40\!\cdot\!0.48/(0.40\!+\!0.48)\!\approx\!0.22\!\geq\!\tau$
and is the sole mask target on this file.

\subsubsection{Multi-Function Group Selection}
\label{sec:method:groups}

Real-world code patches frequently span multiple related
functions~\citep{swebench}, motivating an extension that masks
\emph{groups} of $k=2$ or $k=3$ structurally connected functions; we use
this variant only in the granularity ablation
(Section~\ref{sec:exp:ablation}). A group score $\mathrm{FIM}(G)$
multiplies a coupling term, the harmonic-mean-like product of
group-level $\hat{H}(G)$ and $\hat{I}(G)$, and a difficulty penalty,
with $\hat{I}(G)$ recomputed under joint masking so that intra-group
references cannot inflate the score. Topology patterns (eight in total:
caller-callee, co-callee, sibling-coupled, mutual-call, call-chain, hub,
fan-in, class-triad), the full equations, and selection thresholds are
given in Appendix~\ref{app:algo:groups}.

\subsection{Chain-of-Thought Augmentation}
\label{sec:method:cot}

For each selected target we generate a natural-language rationale via
Gemini-3-Preview prompted with the masked file and the original function
body (template in Appendix~\ref{app:algo:cot}). At training time, the
rationale is embedded \emph{inside} the FIM middle span, immediately
before the function body:
%
\begin{center}\small
\texttt{<fim\_prefix>}\,$\langle$prefix$\rangle$\,%
\texttt{<fim\_suffix>}\,$\langle$suffix$\rangle$\,%
\texttt{<fim\_middle>}\,$\langle$rationale$\rangle$\,$\langle$body$\rangle$
\end{center}
%
The model is therefore trained to produce, in order, a reasoning trace
followed by code consistent with that trace---mirroring the
think-then-act structure of an agent step. The role of the rationale is
to provide a structured intermediate supervision aligned with the FIM
objective rather than to distill the teacher's generic capability; we
isolate this in Section~\ref{sec:exp:ablation} via a \emph{self-CoT}
variant in which the model under training generates its own rationales,
and find that the FIM structure alone already accounts for a substantial
fraction of the gain.
